% !TeX root = ../教材题库_全部条目_转写工作稿.tex
% Source: XJ-B1-C06-EX-P240-Y0489
% 第6章 / 6.4 用样本估计总体 / 教材 p.240、241 / PDF 第247、248页
% OCR未校对：公式与文字均可能有误；未完成数学TeX转写；禁止直接发给学生
% 原题及相邻原解截图（核对用）：../题目原图/XJ-B1-C06-EX-P240-Y0489-01.webp, ../题目原图/XJ-B1-C06-EX-P240-Y0489-02.webp
% 本文件只在转写完成后作为题干使用；OCR 条目尚不是可靠题干。
\begin{bankitem}{6.4 例题 例题5}{第6章 / 6.4 用样本估计总体 / 教材 p.240、241 / PDF 第247、248页}{XJ-B1-C06-EX-P240-Y0489}{OCR未校对：公式与文字均可能有误；未完成数学TeX转写；禁止直接发给学生}
\begin{ocrdraft}
[教材p.240 / PDF第247页]
\par
0〔0]某创业公司全体职工某月的奖金如下：
\par
奖金/元
\par
18000
\par
12000
\par
80006000
\par
4000
\par
2500200015001200
\par
1
\par
2
\par
3
\par
4
\par
20
\par
22
\par
12
\par
6
\par
（总经理）（副总经理）
\par
（1)试求出该公司该月奖金数据的众数、中位数和平均数．
\par
（2）你认为用平均数、中位数或众数中的哪一个更能反映该公司的月奖金水平？
\par
（3）能用该公司的月奖金数据估计全国该类创业公司的月奖金水平吗？说说你
\par
的理由．
\par
解（1）在上述80个数据中，2佣0出现了22次，出现的次数最多，因此这
\par
组数据的众数是20佣．
\par
把这80个数据按从小到大的顺序排列后，位于中间的数是2000，2500，因此
\par
2000十2500
\par
这组数据的中位数是
\par
=2250．
\par
2
\par
这组数据的平均数为
\par
18000十12000\ensuremath{\times}2十8000\ensuremath{\times}3十“，十1200\ensuremath{\times}6一249200
\par
=3115．
\par
若把这组数据的众数、中位数、平均数用图表示，则可得图6，4一1.
\par
中位数平均数
\par
众数《
\par
6人12人22人20人
\par
10人
\par
4人
\par
3人
\par
2人
\par
12001500《《2500
\par
4000
\par
6000
\par
8000
\par
12000
\par
18000
\par
2000
\par
2250
\par
3115
\par
图6．4一1
\par
（2）由于大多数员工的月奖金达不到平均数3115，显然用平均数作为该公司员
\par
[教材p.241 / PDF第248页]
\par
工月奖金的代表值并不合适；众数2000及中位数2250在一定程度上代表了大多数
\par
人的月奖金水平，较能反映月奖金水平的实际情况．
\par
（3）由于全国各地的月奖金差异性较大，因而不能用一个公司的数据估计全国
\par
该类公司的月奖金水平．
\par
平均数、中位数和众数都是一组数据的代表，它们从不同侧面反映了数据的集
\par
中趋势．平均数的计算要用到所有的数据，它能够充分利用数据提供的信息，因此
\par
在现实生活中应用较广，但它容易受极端值的影响；中位数对极端值不敏感，但没
\par
有利用数据中的所有信息；众数只能反映一组数据中出现次数最多的数据，也没有
\par
利用数据中的所有信息．
\end{ocrdraft}
\end{bankitem}
